% GENERATED FILE. Edit canonical lessons, then run npm run generate:books.
% canonical-source-hash: fnv1a64:f1f28d08d4e1b30f
% canonical-lessons: MR-C394-bhan, MR-C394-janiv, MR-C394-jagarukta, MR-C394-vivek, MR-C394-umagne

\chapter{Being Aware}
\label{ch:mr-being-aware}
\hlchaptermodality{394}

\begin{chapteropening}
\textbf{By the end of this chapter:} \emph{I can talk about awareness, realization and good judgement.}

Complete the last lesson of chapter 394: I can talk about awareness, realization and good judgement.
\end{chapteropening}

\section[\texorpdfstring{\mr{भान} (bhān)}{भान (bhān)}]{\mr{भान} (bhān) — awareness; consciousness}

\label{lesson:MR-C394-bhan}



\begin{quote}
Before the new one: say the Marathi for to request, then the Marathi for to follow up.
\end{quote}

\subsection*{You'll want to know: \mr{भान}}
\textbf{\mr{भान}} — \emph{bhān} — \textquotedblleft{}awareness; consciousness\textquotedblright{}.

\begin{quote}
\mr{खेळताना} \mr{मुलांना} \mr{वेळेचं} \mr{भान} \mr{राहिलं} \mr{नाही}.
\end{quote}

\begin{quote}
\emph{kheḷtānā mulānnā veḷecaṁ bhān rāhilaṁ nāhī.}
\end{quote}

\begin{quote}
While playing, the children lost track of time.
\end{quote}

Neuter.

X \mr{चं} \mr{भान} \mr{ठेवणे}, to keep X in mind.

\begin{cousinweb}
\textbf{\mr{भान}} (\emph{bhān}) is a tatsama, taken from Sanskrit \emph{bhāna}, \textquotedblleft{}appearance\textquotedblright{} (CDIAL 9451): what appears to the mind.
\end{cousinweb}

\begin{grammarlens}[title={What you've built}]
One more word for saying what happens if.
\end{grammarlens}

\subsection*{Your turn}
\begin{itemize}
\raggedright
  \item \emph{Say it:} \emph{bhān}
  \item \emph{Say it:} \emph{bhān}, once more
  \item \emph{Say it:} \emph{kheḷtānā mulānnā veḷecaṁ bhān rāhilaṁ nāhī.} — While playing, the children lost track of time.
  \item \emph{Recall:} say the Marathi for to request, then the Marathi for to follow up, then say \emph{bhān} again
\end{itemize}

\begin{tcolorbox}[breakable,colback=teal!4,colframe=teal!35!black,title={Before you move on}]
What does \mr{भान} mean? (\textquotedblleft{}Awareness; consciousness\textquotedblright{}.) Say it once more.
\end{tcolorbox}
\section[\texorpdfstring{\mr{जाणीव} (jāṇīv)}{जाणीव (jāṇīv)}]{\mr{जाणीव} (jāṇīv) — a realization; awareness}

\label{lesson:MR-C394-janiv}



\begin{quote}
Before the new one: say the Marathi for to follow up, then the Marathi for awareness.
\end{quote}

\subsection*{You'll want to know: \mr{जाणीव}}
\textbf{\mr{जाणीव}} — \emph{jāṇīv} — \textquotedblleft{}a realization; awareness\textquotedblright{}.

\begin{quote}
\mr{मला} \mr{माझ्या} \mr{चुकीची} \mr{जाणीव} \mr{झाली}.
\end{quote}

\begin{quote}
\emph{malā mājhyā cukīcī jāṇīv jhālī.}
\end{quote}

\begin{quote}
I realized my mistake.
\end{quote}

Feminine.

X \mr{ची} \mr{जाणीव} \mr{होणे}, to become aware of X.

\begin{cousinweb}
\textbf{\mr{जाणीव}} (\emph{jāṇīv}) is a noun built on \mr{जाणणे}, \textquotedblleft{}to know\textquotedblright{}, with the ending -\mr{ईव}; CDIAL lists the verb under Sanskrit \emph{jānāti}, \textquotedblleft{}knows\textquotedblright{} (CDIAL 5193).
\end{cousinweb}

\begin{grammarlens}[title={What you've built}]
One more word for saying what happens if.
\end{grammarlens}

\subsection*{Your turn}
\begin{itemize}
\raggedright
  \item \emph{Say it:} \emph{jāṇīv}
  \item \emph{Say it:} \emph{jāṇīv}, once more
  \item \emph{Say it:} \emph{malā mājhyā cukīcī jāṇīv jhālī.} — I realized my mistake.
  \item \emph{Recall:} say the Marathi for to follow up, then the Marathi for awareness, then say \emph{jāṇīv} again
\end{itemize}

\begin{tcolorbox}[breakable,colback=teal!4,colframe=teal!35!black,title={Before you move on}]
What does \mr{जाणीव} mean? (\textquotedblleft{}A realization; awareness\textquotedblright{}.) Say it once more.
\end{tcolorbox}
\section[\texorpdfstring{\mr{जागरूकता} (jāgarūktā)}{जागरूकता (jāgarūktā)}]{\mr{जागरूकता} (jāgarūktā) — alertness; public awareness}

\label{lesson:MR-C394-jagarukta}



\begin{quote}
Before the new one: say the Marathi for awareness, then the Marathi for a realization.
\end{quote}

\subsection*{You'll want to know: \mr{जागरूकता}}
\textbf{\mr{जागरूकता}} — \emph{jāgarūktā} — \textquotedblleft{}alertness; public awareness\textquotedblright{}.

\begin{quote}
\mr{आरोग्याबद्दल} \mr{लोकांमध्ये} \mr{जागरूकता} \mr{वाढली} \mr{आहे}.
\end{quote}

\begin{quote}
\emph{ārogyābaddal lokāmmadhye jāgarūktā vāḍhlī āhe.}
\end{quote}

\begin{quote}
Awareness about health has grown among people.
\end{quote}

Feminine.

\mr{जागरूक}, alert, + -\mr{ता}.

\begin{cousinweb}
\textbf{\mr{जागरूकता}} (\emph{jāgarūktā}) is a learned noun built from Sanskrit \emph{jāgarūka}, \textquotedblleft{}wakeful, watchful\textquotedblright{} (Monier-Williams), with \emph{-tā}, \textquotedblleft{}-ness\textquotedblright{}.
\end{cousinweb}

\begin{grammarlens}[title={What you've built}]
One more word for saying what happens if.
\end{grammarlens}

\subsection*{Your turn}
\begin{itemize}
\raggedright
  \item \emph{Say it:} \emph{jāgarūktā}
  \item \emph{Say it:} \emph{jāgarūktā}, once more
  \item \emph{Say it:} \emph{ārogyābaddal lokāmmadhye jāgarūktā vāḍhlī āhe.} — Awareness about health has grown among people.
  \item \emph{Recall:} say the Marathi for awareness, then the Marathi for a realization, then say \emph{jāgarūktā} again
\end{itemize}

\begin{tcolorbox}[breakable,colback=teal!4,colframe=teal!35!black,title={Before you move on}]
What does \mr{जागरूकता} mean? (\textquotedblleft{}Alertness; public awareness\textquotedblright{}.) Say it once more.
\end{tcolorbox}
\section[\texorpdfstring{\mr{विवेक} (vivek)}{विवेक (vivek)}]{\mr{विवेक} (vivek) — good judgement; discretion}

\label{lesson:MR-C394-vivek}



\begin{quote}
Before the new one: say the Marathi for a realization, then the Marathi for alertness.
\end{quote}

\subsection*{You'll want to know: \mr{विवेक}}
\textbf{\mr{विवेक}} — \emph{vivek} — \textquotedblleft{}good judgement; discretion\textquotedblright{}.

\begin{quote}
\mr{रागात} \mr{माणूस} \mr{विवेक} \mr{विसरतो}.
\end{quote}

\begin{quote}
\emph{rāgāt māṇūs vivek visarto.}
\end{quote}

\begin{quote}
In anger a person forgets good judgement.
\end{quote}

Masculine.

\mr{विवेकी}, sensible.

\begin{cousinweb}
\textbf{\mr{विवेक}} (\emph{vivek}) is a tatsama, taken from Sanskrit \emph{viveka}, \textquotedblleft{}discrimination, judgement\textquotedblright{} (Monier-Williams).
\end{cousinweb}

\begin{grammarlens}[title={What you've built}]
One more word for saying what happens if.
\end{grammarlens}

\subsection*{Your turn}
\begin{itemize}
\raggedright
  \item \emph{Say it:} \emph{vivek}
  \item \emph{Say it:} \emph{vivek}, once more
  \item \emph{Say it:} \emph{rāgāt māṇūs vivek visarto.} — In anger a person forgets good judgement.
  \item \emph{Recall:} say the Marathi for a realization, then the Marathi for alertness, then say \emph{vivek} again
\end{itemize}

\begin{tcolorbox}[breakable,colback=teal!4,colframe=teal!35!black,title={Before you move on}]
What does \mr{विवेक} mean? (\textquotedblleft{}Good judgement; discretion\textquotedblright{}.) Say it once more.
\end{tcolorbox}
\section[\texorpdfstring{\mr{उमगणे} (umagṇe)}{उमगणे (umagṇe)}]{\mr{उमगणे} (umagṇe) — to dawn on; to become clear}

\label{lesson:MR-C394-umagne}



\begin{quote}
Before the new one: say the Marathi for alertness, then the Marathi for good judgement.
\end{quote}

\subsection*{You'll want to know: \mr{उमगणे}}
\textbf{\mr{उमगणे}} — \emph{umagṇe} — \textquotedblleft{}to dawn on; to become clear\textquotedblright{}.

\begin{quote}
\mr{खूप} \mr{वेळाने} \mr{मला} \mr{त्याच्या} \mr{बोलण्याचा} \mr{अर्थ} \mr{उमगला}.
\end{quote}

\begin{quote}
\emph{khūp veḷāne malā tyācyā bolaṇyācā artha umaglā.}
\end{quote}

\begin{quote}
Much later the meaning of his words dawned on me.
\end{quote}

Past (to him …): \mr{त्याला} \mr{उमगलं}, \emph{tyālā umaglaṁ}.

Dative: \mr{मला} \mr{उमगलं}, it dawned on me.

\begin{cousinweb}
\textbf{\mr{उमगणे}} (\emph{umagṇe}) is inherited: CDIAL lists Marathi \emph{umagṇẽ}, \textquotedblleft{}to search out; to become known\textquotedblright{}, under \emph{*unmārgati}, \textquotedblleft{}searches out\textquotedblright{} (CDIAL 2126).
\end{cousinweb}

\begin{grammarlens}[title={What you've built}]
One more word for saying what happens if.
\end{grammarlens}

\subsection*{Your turn}
\begin{itemize}
\raggedright
  \item \emph{Say it:} \emph{umagṇe}
  \item \emph{Say it:} \emph{umagṇe}, once more
  \item \emph{Say it:} \emph{khūp veḷāne malā tyācyā bolaṇyācā artha umaglā.} — Much later the meaning of his words dawned on me.
  \item \emph{Recall:} say the Marathi for alertness, then the Marathi for good judgement, then say \emph{umagṇe} again
\end{itemize}

\begin{tcolorbox}[breakable,colback=teal!4,colframe=teal!35!black,title={Before you move on}]
What does \mr{उमगणे} mean? (\textquotedblleft{}To dawn on; to become clear\textquotedblright{}.) Say it once more.
\end{tcolorbox}
